% Manual · parte inicial: organización, banco, seguridad e hilo conductor.

\PortadillaLaboratorio{Antes de empezar · 10 minutos de lectura}{Cómo funciona este laboratorio}{En dos sesiones va a construir dos receptores reales. En la primera, una emisora de FM pasa por una antena, un HackRF y un programa hasta convertirse en sonido. En la segunda, ráfagas de microsegundos emitidas por aviones se convierten en una tabla y un mapa. En ambas aprenderá lo mismo: distinguir lo que \emph{ve} de lo que puede \emph{demostrar}.}

\section{Cuatro documentos, una tarea}

\begin{longtable}{@{}>{\raggedright\arraybackslash}p{0.27\linewidth}>{\raggedright\arraybackslash}p{0.66\linewidth}@{}}
\toprule
\textbf{Documento} & \textbf{Para qué lo usa} \\
\midrule
\endhead
\textbf{Manual} (este) & Qué hacer, en qué orden, con qué valores y qué debe ver. Es lo único que necesita abierto junto al equipo. \\
\textbf{Anexo A} · Fundamentos & Por qué funciona: teoría, aclaraciones, seguridad detallada, diagnóstico de fallos y datos de referencia. \\
\textbf{Anexo B} · Evaluación y registro & Cuestionarios, hojas de registro y entregables. Imprímalo y llévelo. \\
\textbf{Guía del instructor} & Preparación del banco, claves y rúbrica. No es para el estudiante. \\
\bottomrule
\end{longtable}

\subsection{Cómo leer un paso}

Cada paso importante es un recuadro con una acción y hasta seis filas:

\Paso{La acción, en una frase}
{El conector, el menú o el control exactos.}
{El valor que debe poner.}
{La razón técnica, en una línea.}
{Lo que verá u oirá si todo va bien.}
{Una comprobación concreta.}
{Qué hacer de inmediato si no sale.}

Las órdenes de terminal (Laboratorio 02) se presentan en otro recuadro: la orden exacta arriba y, debajo, su propósito, la salida esperada, cómo reconocer un fallo y cómo recuperarse. Cuando un paso le pida registrar algo, el manual le indicará la hoja del Anexo B, por ejemplo la \hoja{W2}.

\section{El banco}
\label{sec:banco}

\begin{longtable}{@{}>{\raggedright\arraybackslash}p{0.34\linewidth}>{\raggedright\arraybackslash}p{0.42\linewidth}>{\raggedright\arraybackslash}p{0.16\linewidth}@{}}
\toprule
\textbf{Elemento} & \textbf{Función} & \textbf{Laboratorio} \\
\midrule
\endhead
HackRF One y su cable USB de datos & Receptor SDR de 1 MHz a 6 GHz & 01 y 02 \\
Antena telescópica ANT500 (75 MHz--1 GHz, SMA macho) & Antena de FM & 01 \\
Antena de 1090 MHz con base magnética y 3 m de cable (SMA macho) & Antena de ADS-B & 02 \\
Filtro pasabanda de 1090 MHz, pasivo (SMA) & Atenúa emisiones lejanas a 1090 MHz & 02 \\
Cable de extensión SMA de bajas pérdidas & Llevar la antena a la ventana & 01 y 02 \\
Adaptadores SMA del paquete del HackRF & Ajustar géneros de conector & 01 y 02 \\
Analizador SIGLENT SSA3032X Plus & Medidas de espectro & 01 y 02 \\
Adaptador N macho -- SMA hembra & Conectar las antenas a la entrada N del analizador & 01 y 02 \\
PC con Windows 10/11, SDRangel 7.27.2 y radioconda; auriculares & Estación del Laboratorio 01 & 01 \\
PC con Ubuntu 24.04 LTS, SDRangel 7.27.2 (snap) y \texttt{hackrf} & Estación del Laboratorio 02 & 02 \\
Grabación de referencia \texttt{datos\_adsb/\ldots.sdriq} & Contraste con datos conocidos & 02 \\
\bottomrule
\end{longtable}

\section{Tres reglas que no se negocian}
\label{sec:reglas}

\begin{alerta}[Seguridad]
\begin{enumerate}
  \item \textbf{Solo se recibe.} No añada en SDRangel nada que diga \emph{Output}, \emph{Modulator} o \emph{TX}, y no conecte nunca el HackRF al analizador. Emitir en 1030/1090 MHz interfiere con sistemas de seguridad aérea.
  \item \textbf{Las entradas son frágiles.} Toque una superficie metálica antes de manipular conectores. Conecte y desconecte con el receptor detenido. Apriete las tuercas SMA a mano y no gire el cuerpo de la antena.
  \item \textbf{La alimentación de antena se queda apagada.} Las casillas \gui{Bias tee} y \gui{RF amp} de SDRangel permanecen desactivadas salvo indicación expresa del instructor.
\end{enumerate}
Detalle de límites, descargas electrostáticas y emplazamientos cercanos a radares: \anexoA{sec:a-seguridad}.
\end{alerta}

\section{La idea que une los dos laboratorios}

Una pantalla llena de colores no prueba nada por sí sola. Cada vez que observe algo, pregúntese qué demuestra:

\begin{enumerate}
  \item \textbf{Energía}: hay actividad en una frecuencia.
  \item \textbf{Demodulación}: se recuperó la señal que transportaba la portadora (el audio).
  \item \textbf{Sincronización}: el receptor reconoció la estructura de un mensaje (un preámbulo).
  \item \textbf{Trama válida}: unos bits superaron la comprobación de su protocolo.
  \item \textbf{Información}: campos transmitidos por el emisor (dirección, altitud).
  \item \textbf{Información derivada}: algo que calculó el receptor (una posición).
  \item \textbf{Información externa}: algo que vino de una base de datos (el modelo del avión).
\end{enumerate}

Cada entregable le pedirá declarar el nivel de su evidencia. La escalera completa, con ejemplos, está en el \anexoA{sec:a-escalera}.
